\documentclass[11pt,a4paper]{article}
\usepackage[margin=2cm]{geometry}
\usepackage[utf8]{inputenc}
\usepackage[T1]{fontenc}
\usepackage{booktabs}
\usepackage{xcolor}
\usepackage{parskip}
\usepackage{titlesec}
\usepackage{fancyhdr}
\usepackage{array}
\usepackage{multirow}
\usepackage{mdframed}
\usepackage{enumitem}
\usepackage{hyperref}
\usepackage{lmodern}
\usepackage{microtype}

% ── colours ─────────────────────────────────────────────────────────────────
\definecolor{stageblue}{HTML}{1E3A5F}
\definecolor{accentblue}{HTML}{2563EB}
\definecolor{lightblue}{HTML}{EFF6FF}
\definecolor{green}{HTML}{16A34A}
\definecolor{amber}{HTML}{D97706}
\definecolor{purple}{HTML}{7C3AED}
\definecolor{gray}{HTML}{6B7280}
\definecolor{lightgray}{HTML}{F3F4F6}

% ── header/footer ───────────────────────────────────────────────────────────
\pagestyle{fancy}
\fancyhf{}
\lhead{\textcolor{gray}{\small\textit{Pipeline Scripts Guide --- AI Coding Agent Adoption Study}}}
\rhead{\textcolor{gray}{\small\textit{Manuel Youssef Haik Kevorkian}}}
\cfoot{\textcolor{gray}{\small\thepage}}
\renewcommand{\headrulewidth}{0.3pt}

% ── section styling ─────────────────────────────────────────────────────────
\titleformat{\section}[block]
  {\color{white}\normalfont\large\bfseries}
  {}
  {0pt}
  {\colorbox{stageblue}{\parbox{\dimexpr\linewidth-2\fboxsep}{\thesection\quad #1}}}

\titleformat{\subsection}[block]
  {\color{accentblue}\normalfont\normalsize\bfseries}
  {\thesubsection}
  {6pt}
  {#1}

% ── custom environments ─────────────────────────────────────────────────────
\newmdenv[
  linecolor=accentblue,
  backgroundcolor=lightblue,
  linewidth=1.2pt,
  leftmargin=0pt, rightmargin=0pt,
  innerleftmargin=8pt, innerrightmargin=8pt,
  innertopmargin=6pt, innerbottommargin=6pt,
  skipabove=4pt, skipbelow=4pt
]{infobox}

\newmdenv[
  linecolor=green,
  backgroundcolor=white,
  linewidth=1pt,
  leftmargin=0pt, rightmargin=0pt,
  innerleftmargin=8pt, innerrightmargin=8pt,
  innertopmargin=5pt, innerbottommargin=5pt,
  skipabove=3pt, skipbelow=3pt
]{outputbox}

\newmdenv[
  linecolor=amber,
  backgroundcolor=white,
  linewidth=1pt,
  leftmargin=0pt, rightmargin=0pt,
  innerleftmargin=8pt, innerrightmargin=8pt,
  innertopmargin=5pt, innerbottommargin=5pt,
  skipabove=3pt, skipbelow=3pt
]{numbersbox}

\setlist[itemize]{noitemsep, topsep=2pt, leftmargin=*}
\setlength{\parskip}{4pt}

% ── title ───────────────────────────────────────────────────────────────────
\title{%
  \vspace{-1cm}
  \textbf{\color{stageblue}Pipeline Scripts Guide}\\[0.3em]
  \large Coding Agent Adoption in JavaScript/TypeScript/npm\\
  Repositories on GitHub\\[0.4em]
  \normalsize One Page Per Script --- Inputs, Outputs, Methods, and Thesis Contribution
}
\author{Manuel Youssef Haik Kevorkian}
\date{October 2026 \quad \textcolor{red}{\textbf{DRAFT}}}

\begin{document}
\maketitle
\thispagestyle{fancy}
\vspace{-0.5cm}
\hrule
\vspace{0.3cm}

{\small
This document describes every pipeline script across the four stages of the study.
Each page covers one script: what it does, what it reads, what it produces,
the key numbers it generated, and how it contributes to answering the three
research questions. Scripts are presented in execution order.
The pipeline has \textbf{13 scripts} organised into 4 stages and produces
a final corpus of \textbf{60,129 repositories} with complete adoption labels
and geographic metadata.
}

\vspace{0.4cm}
\begin{center}
\begin{tabular}{@{}clll@{}}
\toprule
\textbf{Stage} & \textbf{Script} & \textbf{Purpose} & \textbf{Key output} \\
\midrule
1a & stage1a\_collect\_repos.py    & GitHub API collection        & 60,129 repos \\
1b & stage1b\_categorize\_ecosystem.py & Lockfile detection      & ecosystem labels \\
2a & stage2a\_location\_fast.py   & Location field lookup        & 13,625 owners located \\
2b & stage2b\_location\_deep.py   & Bio / TLD / network scan     & 8,323 more located \\
3a & stage3a\_config\_files.py    & Config file existence check  & 25,010 agentic \\
3b & stage3b\_commit\_trailers.py & Commit trailer scan          & 9,618 agentic + 3,846 ai-assisted \\
3c & stage3c\_deep\_scan.py       & Full GraphQL deep scan       & 764 agentic + 4,657 ai-assisted \\
3d & stage3d\_reclassify.py       & Reclassification (Pass 4)    & 258 upgraded to agentic \\
4a & stage4a\_geographic\_analysis.py & Country join + ranking   & 65-country table \\
4b & stage4b\_statistical\_tests.py   & RQ1 / RQ2 / RQ3 tests    & 3 statistical tests \\
4c & stage4c\_breakdowns.py           & Ecosystem + tool analysis & breakdown CSVs \\
4d & stage4d\_size\_analysis.py        & Repo-size analysis        & 3-approach size table \\
4e & stage4e\_mena\_analysis.py        & MENA regional deep-dive   & MENA ranking + stats \\
\bottomrule
\end{tabular}
\end{center}

\newpage

% ══════════════════════════════════════════════════════════════════════════════
\section{Stage 1 --- Corpus Construction}
% ══════════════════════════════════════════════════════════════════════════════

\subsection*{\texttt{stage1a\_collect\_repos.py} --- GitHub API Repository Collection}

\begin{infobox}
\textbf{Purpose:} Queries the GitHub GraphQL Search API to collect all JavaScript
and TypeScript repositories that meet the corpus inclusion criteria.
Produces the raw population to which all subsequent analyses are applied.
No sampling is performed --- the output is a complete census.
\end{infobox}

\textbf{Method.}
The script issues two separate search queries (one for JavaScript, one for
TypeScript, because the API does not support combined language filters) using
rolling 5-day creation date windows from 24~February~2025 to October~2026.
Any window that hits the API's 1,000-result cap is recursively split into
2-day sub-windows until every window fits under the cap. For each repository
returned, the script records: full name, stars, forks, commit count,
JavaScript+TypeScript byte count (from GitHub Linguist), URL, primary
language, and creation date. After collection, duplicate entries and
repositories failing any inclusion filter are removed.

\textbf{Inclusion filters:}
\begin{itemize}
  \item Created on or after \textbf{24 February 2025} (start of the agentic era)
  \item \textbf{$\geq 4$ stars} (excludes abandoned toy projects)
  \item \textbf{$\geq 50$ commits} (excludes minimal/tutorial repositories)
  \item \textbf{Not a fork} (only original codebases)
  \item \textbf{JS+TS bytes $\geq 175{,}000$} (ensures meaningful code volume)
\end{itemize}

\begin{numbersbox}
\textbf{Key numbers generated:}\\[2pt]
\begin{tabular}{@{}ll@{}}
Total repos collected after all filters: & \textbf{60,129} \\
TypeScript repos:                        & 45,314 (75.4\%) \\
JavaScript repos:                        & 14,815 (24.6\%) \\
Stars range:                             & 4 -- 391,443 \quad median: 10 \\
Commits range:                           & 50 -- 1,003,888 \quad median: 177 \\
JS+TS bytes median:                      & 837 KB \\
\end{tabular}
\end{numbersbox}

\textbf{Input:} GitHub GraphQL Search API (live, no local input file)

\textbf{Output:} \texttt{data/stage1/global\_corpus.csv} --- 60,129 rows,
13 columns: \texttt{full\_name, stars, forks, commits, lang\_bytes, url,
owner\_login, repo\_name, primary\_language, size\_stratum, created\_at,
ecosystem\_type, corpus\_source}

\begin{outputbox}
\textbf{Thesis contribution.}
This script defines the study population. The 24~Feb~2025 threshold ensures
every repository in the corpus was created during the period when developers
had simultaneous access to mature AI coding agents (Claude Code, GitHub
Copilot, Cursor, Aider, etc.). Including pre-threshold repositories would
mix projects that never had realistic access to agents with those that did,
artificially suppressing the adoption rate. The 60,129 repositories
constitute the denominator for all subsequent adoption rate calculations,
including RQ1's 73.0\% global adoption figure.
\end{outputbox}

\newpage

\subsection*{\texttt{stage1b\_categorize\_ecosystem.py} --- Package Manager Classification}

\begin{infobox}
\textbf{Purpose:} Classifies each of the 60,129 repositories into a package
manager ecosystem by detecting lockfiles and configuration files at the
repository root via the GitHub GraphQL API.
\end{infobox}

\textbf{Method.}
For each repository, a single GraphQL query checks for the presence of known
lockfiles in order of priority: \texttt{bun.lockb} (Bun), \texttt{deno.lock}
(Deno), \texttt{pnpm-lock.yaml} (pnpm), \texttt{yarn.lock} (Yarn),
\texttt{package-lock.json} (npm). If a \texttt{package.json} is present
without any lockfile the repo is labelled \texttt{npm\_no\_lock}.
If no \texttt{package.json} is found the repo is labelled \texttt{vanilla}
(a JavaScript project not using the npm registry). The priority order
prevents double-counting in polyglot projects that commit multiple lockfiles.

\begin{numbersbox}
\textbf{Key numbers generated:}\\[2pt]
\begin{tabular}{@{}lrr@{}}
\toprule
Ecosystem      & Repositories & Share \\
\midrule
npm            & 28,248 & 47.0\% \\
pnpm           & 12,229 & 20.3\% \\
vanilla        & 10,177 & 16.9\% \\
npm\_no\_lock  &  7,979 & 13.3\% \\
yarn           &  1,091 &  1.8\% \\
deno           &    250 &  0.4\% \\
bun            &    107 &  0.2\% \\
unknown        &     48 &  0.1\% \\
\midrule
\textbf{Total} & \textbf{60,129} & \textbf{100\%} \\
\bottomrule
\end{tabular}
\end{numbersbox}

\textbf{Input:} \texttt{data/stage1/global\_corpus.csv} (60,129 repos)

\textbf{Output:} Ecosystem labels written back into \texttt{global\_corpus.csv}
as the \texttt{ecosystem\_type} column. Also feeds \texttt{stage4c\_breakdowns.py}
which computes per-ecosystem adoption rates.

\begin{outputbox}
\textbf{Thesis contribution.}
Ecosystem classification reveals that adoption is not uniform across
package managers. Stage 4c shows pnpm leads at 82.5\%, followed by npm at
73.7\% and vanilla at 57.5\%. This variation is discussed in the Results
chapter as evidence that more engineering-focused package manager choices
(pnpm requires deliberate configuration) correlate with higher rates of
AI coding agent adoption. The ecosystem label also enables the Results chapter
to contextualise why the corpus contains disproportionately many pnpm
repositories relative to global market share.
\end{outputbox}

\newpage

% ══════════════════════════════════════════════════════════════════════════════
\section{Stage 2 --- Owner Location Inference}
% ══════════════════════════════════════════════════════════════════════════════

\subsection*{\texttt{stage2a\_location\_fast.py} --- Pass 1: Location Field Lookup}

\begin{infobox}
\textbf{Purpose:} Infers the country of each repository owner by reading
the \texttt{location} field of their GitHub profile. This is the cheapest
and most accurate signal, processed in a single batched pass.
\end{infobox}

\textbf{Method.}
The 60,129 repositories have \textbf{50,002 unique owners}. For each owner,
the script queries the GitHub REST API for the profile \texttt{location}
string. The raw string is normalised against a country name and alias
dictionary (covering abbreviations, demonyms, and common misspellings)
to produce a standardised country name. Owners whose location string is
empty, clearly fake (e.g.\ ``Earth'', ``Internet''), or ambiguous
(e.g.\ city name shared by multiple countries) are marked
\texttt{not\_specified} and passed to Stage~2b.

\begin{numbersbox}
\textbf{Key numbers generated:}\\[2pt]
\begin{tabular}{@{}ll@{}}
Unique owners queried:           & 50,002 \\
Located via location field:      & \textbf{13,625} (27.2\% of owners) \\
Marked not\_specified:           & 36,377 (72.8\%) passed to Stage 2b \\
Top countries (Pass 1):          & China 1,521 \quad US 1,491 \quad India 1,650 \\
\end{tabular}
\end{numbersbox}

\textbf{Input:} \texttt{data/stage1/global\_corpus.csv} (owner\_login column)

\textbf{Output:} \texttt{data/stage2/owner\_locations\_fast.csv} ---
50,002 rows with columns: \texttt{owner\_login, location\_raw,
inferred\_country, confidence, source}

\begin{outputbox}
\textbf{Thesis contribution.}
The location field is the only self-reported, unambiguous location signal
on GitHub. It locates 27.2\% of owners with high confidence in a single
cheap pass. These are the most reliable country assignments in the study
and form the foundation of the geographic ranking in Stage~4a.
The 72.8\% that remain unresolved motivate the design of Stage~2b and
directly lead to the study's key limitation: 56.1\% of owners cannot
be assigned a country, which is reported as a validity threat in the
Methodology chapter.
\end{outputbox}

\newpage

\subsection*{\texttt{stage2b\_location\_deep.py} --- Pass 2: Bio, TLD, and Network Inference}

\begin{infobox}
\textbf{Purpose:} Attempts to infer the country of the 36,377 owners that
Pass~1 could not resolve, using six additional signals extracted from the
GitHub profile and the owner's social graph.
\end{infobox}

\textbf{Method.}
Six inference signals are tried in priority order:
\begin{itemize}
  \item \textbf{Bio text} --- scans the profile biography for country names,
        city names, or language-specific phrases.
  \item \textbf{Company name} --- maps company names to headquarters country.
  \item \textbf{Email TLD} --- e.g.\ \texttt{.de} infers Germany.
  \item \textbf{Website TLD} --- same logic applied to the profile website.
  \item \textbf{Organisation location} --- if the owner belongs to a GitHub
        organisation, the organisation's location field is used.
  \item \textbf{Network majority} --- follows up to 20 mutual followers
        and takes the majority country of those already resolved.
\end{itemize}
A confidence score (low / medium / high) is assigned based on which signal
resolved the owner and whether signals agreed.

\begin{numbersbox}
\textbf{Key numbers generated:}\\[2pt]
\begin{tabular}{@{}ll@{}}
Input (unresolved from Pass 1):  & 36,377 owners \\
Located by bio text:             & 232 \\
Located by company name:         & 144 \\
Located by email TLD:            & 263 \\
Located by website TLD:          & 439 \\
Located by org location:         & 892 \\
Located by network majority:     & 6,333 \\
\textbf{Total resolved (Pass 2):}& \textbf{8,323} owners \\
Still not specified after Pass 2:& 28,054 (56.1\% of all owners) \\
\end{tabular}
\end{numbersbox}

\textbf{Input:} Owners marked \texttt{not\_specified} in
\texttt{owner\_locations\_fast.csv} (36,377 rows)

\textbf{Output:} \texttt{data/stage2/owner\_locations\_deep.csv} and
\texttt{data/stage2/owner\_locations\_all.csv} (merged Pass~1 + Pass~2,
50,002 rows) and \texttt{data/stage2/not\_specified\_owners.csv} (28,054 rows)

\begin{outputbox}
\textbf{Thesis contribution.}
Pass~2 increases the located population from 13,625 to 21,948 owners
(+61\%), more than doubling geographic coverage. The dominant signal
is network majority (6,333 recoveries), which exploits the geographic
clustering of developer social graphs. Together, Passes~1 and~2 enable
Stage~4a to build a 65-country ranking covering 26,586 repositories
(44.2\% of the corpus). The 56.1\% that remain unresolved is reported
transparently as a key limitation: if unresolved owners differ
systematically from resolved ones (e.g.\ lower-income developers
less likely to complete profiles), country-level rates will be biased upward
for higher-income countries.
\end{outputbox}

\newpage

% ══════════════════════════════════════════════════════════════════════════════
\section{Stage 3 --- AI Adoption Detection}
% ══════════════════════════════════════════════════════════════════════════════

\subsection*{\texttt{stage3a\_config\_files.py} --- Pass 1: Configuration File Detection}

\begin{infobox}
\textbf{Purpose:} Detects AI coding agent configuration files at the
repository root for all 60,129 repositories. This is the cheapest signal
in the pipeline: file existence only, no content reading. Repositories
classified here are not re-examined by any later pass.
\end{infobox}

\textbf{Method.}
A single batched GraphQL query per repository checks whether any of the
known agent config filenames exist at the root. The checked filenames are:
\texttt{CLAUDE.md}, \texttt{AGENTS.md}, \texttt{.github/copilot-instructions.md},
\texttt{.cursorrules}, \texttt{.cursorignore}, \texttt{GEMINI.md},
\texttt{.windsurfrules}, \texttt{.roomodes}, and several others.
A positive match immediately classifies the repository as
\textbf{agentic} under the MEDIUM signal family (\textit{config\_file}).
Repositories with no match are written to the unclassified list for Pass~2.

\begin{numbersbox}
\textbf{Key numbers generated:}\\[2pt]
\begin{tabular}{@{}ll@{}}
Input:                               & 60,129 repositories \\
Classified agentic:                  & \textbf{25,010} (41.6\%) \\
Passed to Pass 2 (unclassified):     & 35,119 (58.4\%) \\[4pt]
Top config files found: & \\
\quad CLAUDE.md                      & 15,152 \quad (60.6\% of Pass~1 agentic) \\
\quad AGENTS.md                      & 8,705 \quad (34.8\%) \\
\quad .github/copilot-instructions.md & 715 \quad (2.9\%) \\
\quad .cursorrules                   & 177 \quad (0.7\%) \\
\quad GEMINI.md                      & 106 \quad (0.4\%) \\
\end{tabular}
\end{numbersbox}

\textbf{Input:} \texttt{data/stage1/global\_corpus.csv}

\textbf{Output:}
\texttt{data/stage3/pass1\_results.csv} (60,129 rows) and
\texttt{data/stage3/pass1\_unclassified.csv} (35,119 rows for Pass~2)

\begin{outputbox}
\textbf{Thesis contribution.}
Pass~1 single-handedly detects 41.6\% of the corpus as agentic and accounts
for the majority of all agentic classifications. The dominance of
\texttt{CLAUDE.md} (60.6\% of Pass~1 hits) is one of the most striking
findings in the thesis: Claude Code is by far the most widely adopted
agent in the JS/TS/npm ecosystem. This pass is also central to RQ2 ---
comparing Pass~1-only detection (41.6\%) with the full four-pass pipeline
(73.0\%) reveals the \textbf{+31.4~percentage-point undercounting gap}
that single-signal methods produce.
\end{outputbox}

\newpage

\subsection*{\texttt{stage3b\_commit\_trailers.py} --- Pass 2: Commit Trailer and PR Branch Scan}

\begin{infobox}
\textbf{Purpose:} Scans the commit history and pull request branch names
of the 35,119 repositories that Pass~1 did not classify.
Commit trailers are STRONG signals --- they confirm an agent
\emph{performed} work, not merely that it was configured.
\end{infobox}

\textbf{Method.}
For each unclassified repository, the script fetches the most recent
\textbf{100 commits} via the GitHub REST API and scans each commit message
for agent-specific trailer patterns using regular expressions:
\begin{itemize}
  \item \texttt{Co-authored-by:} lines containing Claude Code, Copilot,
        OpenAI Codex, Cursor, Gemini, Devin, etc.
  \item \texttt{Generated by:} / \texttt{Committed by:} patterns
  \item PR branch names prefixed with \texttt{copilot/}, \texttt{codex/},
        \texttt{devin/}, \texttt{claude/}, \texttt{jules-}
\end{itemize}
Classification rule: $\geq 3$ AI commits or trailers from 2+ distinct agents
$\rightarrow$ \textbf{agentic}. 1--2 AI commits from 1 agent $\rightarrow$
\textbf{ai\_assisted}. No match $\rightarrow$ passed to Pass~3.

\begin{numbersbox}
\textbf{Key numbers generated:}\\[2pt]
\begin{tabular}{@{}ll@{}}
Input:                           & 35,119 repositories \\
Classified agentic:              & \textbf{9,618} (27.4\% of input) \\
Classified ai-assisted:          & \textbf{3,846} (10.9\%) \\
Passed to Pass 3:                & 21,655 (61.7\%) \\[4pt]
Top agents detected in trailers: & \\
\quad Claude Code                & 8,224 repos \\
\quad GitHub Copilot             & 1,637 repos \\
\quad OpenAI Codex               & 1,577 repos \\
\quad Cursor                     & 775 repos \\
\quad Claude Code + Cursor (pair)& 517 repos \\
\end{tabular}
\end{numbersbox}

\textbf{Input:} \texttt{data/stage3/pass1\_unclassified.csv} (35,119 rows)

\textbf{Output:} \texttt{data/stage3/pass2\_results.csv} and
\texttt{data/stage3/pass2\_unclassified.csv} (21,655 rows for Pass~3)

\begin{outputbox}
\textbf{Thesis contribution.}
Pass~2 adds 9,618 agentic repositories that would have been missed by a
config-file-only approach. Commit trailers are STRONG signals because they
prove the agent \emph{authored} commits --- the agent did not just sit in a
config file. The 3,846 ai-assisted repos from this pass all had 1--2 trailer
hits, which Stage~3d (Pass~4) later re-examines; ultimately \textbf{all
3,846} are upgraded to agentic because a single commit trailer is a STRONG
signal. Pass~2 also provides the most granular per-agent breakdown in the
corpus, directly enabling the tool distribution analysis in Stage~4c.
\end{outputbox}

\newpage

\subsection*{\texttt{stage3c\_deep\_scan.py} --- Pass 3: Deep GraphQL Scan}

\begin{infobox}
\textbf{Purpose:} Performs the most thorough examination of the 21,655
repositories still unclassified after Passes~1 and~2. Uses batched
GraphQL queries to fetch commit history, pull requests, CI workflows, and
\texttt{.gitignore} entries, then applies the full three-tier signal hierarchy.
\end{infobox}

\textbf{Method.}
Repositories are batched in groups of 5 per GraphQL query (4--5$\times$ faster
than individual REST calls). For each repository the query fetches:
\textbf{150 commits} and \textbf{50 pull requests}. The signal families
checked are:
\begin{itemize}
  \item \textbf{STRONG}: \textit{agent\_attribution} (commit trailers / author email),
        \textit{pr\_branch} (agent-prefixed PR branches),
        \textit{bot\_author} (known agent bot accounts as commit authors)
  \item \textbf{MEDIUM}: \textit{pr\_label} (PR label containing agent name),
        \textit{tool\_configuration} (\texttt{.gitignore} entries for agentic tools)
  \item \textbf{WEAK}: \textit{pr\_body} (agent mention in PR description),
        \textit{gitignore\_mixed} (editor-with-agent-mode \texttt{.gitignore} entries)
\end{itemize}
Classification rule: $\geq 1$ STRONG $\rightarrow$ \textbf{agentic};
$\geq 2$ distinct MEDIUM families $\rightarrow$ \textbf{agentic};
1 MEDIUM only $\rightarrow$ \textbf{ai\_assisted};
$\geq 1$ WEAK $\rightarrow$ \textbf{ai\_assisted}; nothing $\rightarrow$ \textbf{none}.

\begin{numbersbox}
\textbf{Key numbers generated:}\\[2pt]
\begin{tabular}{@{}ll@{}}
Input:                           & 21,655 repositories \\
Classified agentic:              & \textbf{764} (3.5\% of input) \\
Classified ai-assisted:          & \textbf{4,657} (21.5\%) \\
Classified none:                 & 16,234 (75.0\%) --- final none label \\[4pt]
Top STRONG signals (agentic):    & \\
\quad pr\_branch                 & 704 repos \\
\quad bot\_author                & 45 repos \\
\end{tabular}
\end{numbersbox}

\textbf{Input:} \texttt{data/stage3/pass2\_unclassified.csv} (21,655 rows)

\textbf{Output:} \texttt{data/stage3/pass3\_results.csv} and
\texttt{data/stage3/corpus\_detected\_full.csv} (merged Passes~1--3)

\begin{outputbox}
\textbf{Thesis contribution.}
Pass~3 is the most computationally expensive pass ($\sim$3--4 hours runtime)
but essential for completeness. Its 764 agentic and 4,657 ai-assisted
classifications would all be missed by single-signal approaches. The 16,234
\textbf{none} labels produced here are the definitive ``no evidence found''
classification --- these repos passed through three increasingly thorough
checks and showed no AI signal. The 4,657 ai-assisted repos from this pass
enter Stage~3d (Pass~4) for a final deep re-examination.
\end{outputbox}

\newpage

\subsection*{\texttt{stage3d\_reclassify.py} --- Pass 4: Reclassification of AI-Assisted Repos}

\begin{infobox}
\textbf{Purpose:} Takes every \textbf{ai-assisted} repository produced by
Passes~2 and~3 (8,503 total) and applies two phases of deeper analysis
to upgrade as many as possible to \textbf{agentic}. Pass~4 is the final
classification step; its output is the definitive adoption label for all
60,129 repositories.
\end{infobox}

\textbf{Phase 1 (no API calls).}
Re-applies the tiered classification rule to the 3,846 ai-assisted repos
from Pass~2. Under the strict rule, a single commit trailer is a STRONG
(\textit{agent\_attribution}) signal $\rightarrow$ immediately agentic.
Pass~2 had used a looser count-based rule; Phase~1 corrects this without
any additional API calls. Result: \textbf{all 3,846} upgraded to agentic.

\textbf{Phase 2 (deep API checks).}
Runs five new check families on the 4,657 ai-assisted repos from Pass~3:
\begin{itemize}
  \item Extended config files in subdirectories (e.g.\ \texttt{.claude/settings.json})
  \item Bot account authorship in up to 300 commits
  \item Agent-style PR branch prefixes (expanded set)
  \item CI workflow files (\texttt{.github/workflows/*.yml}) containing agent actions
  \item PR labels and PR body patterns (extended regex)
\end{itemize}

\begin{numbersbox}
\textbf{Key numbers generated:}\\[2pt]
\begin{tabular}{@{}ll@{}}
Input (total ai-assisted):              & 8,503 repositories \\
Phase 1 --- upgraded (Pass~2 repos):   & \textbf{3,846} $\rightarrow$ agentic (100\%) \\
Phase 2 --- input:                     & 4,657 repositories (Pass~3 ai-assisted) \\
Phase 2 --- upgraded to agentic:       & \textbf{258} \\
Phase 2 --- remain ai-assisted:        & 4,399 \\[4pt]
Top new signals (Phase 2):             & \\
\quad \texttt{.mcp.json} (MCP tool)    & 63 repos \\
\quad \texttt{.claude/settings.json}   & 35 repos \\
\quad \texttt{.claude/settings.local.json} & 30 repos \\
\quad Copilot CI workflow              & 13 repos \\
\quad Jules bot author                 & 10 repos \\
\end{tabular}
\end{numbersbox}

\textbf{Input:} \texttt{data/stage3/corpus\_detected\_full.csv} (all ai-assisted rows)

\textbf{Output:} \texttt{data/stage3/reclassification\_results.csv} and
\texttt{data/stage3/corpus\_detected\_full\_v2.csv} --- the final definitive
labels for all 60,129 repositories

\begin{outputbox}
\textbf{Thesis contribution.}
Pass~4 adds the last \textbf{4,104 agentic} repositories (3,846 + 258)
that no prior pass could confirm. Its output is the ground truth used in
every table, chart, and statistical test in the thesis. The final
classification is: 39,496 agentic (65.7\%), 4,399 ai-assisted (7.3\%),
16,234 none (27.0\%). Combined, 43,895 repositories (73.0\%) show evidence
of AI coding agent adoption --- the answer to \textbf{RQ1}. The 4,399
ai-assisted repos represent borderline evidence: configured but unconfirmed.
They are included in the 73.0\% total; excluding them gives the conservative
65.7\% agentic-only figure reported alongside.
\end{outputbox}

\newpage

% ══════════════════════════════════════════════════════════════════════════════
\section{Stage 4 --- Geographic and Statistical Analysis}
% ══════════════════════════════════════════════════════════════════════════════

\subsection*{\texttt{stage4a\_geographic\_analysis.py} --- Country Join and Ranking}

\begin{infobox}
\textbf{Purpose:} Joins the final adoption labels from Stage~3 with the
owner locations from Stage~2 to produce a per-country adoption ranking.
Answers the geographic dimension of \textbf{RQ3}.
\end{infobox}

\textbf{Method.}
Each repository's owner is looked up in \texttt{owner\_locations\_all.csv};
if a country is found, the repository inherits that country label.
Countries with fewer than 30 repositories are excluded from the ranking
(the MIN\_REPOS threshold) to ensure statistically meaningful confidence
intervals. For each eligible country, the script computes: total repos,
agentic count, ai-assisted count, adoption rate, and Wilson 95\% confidence
interval. Countries are ranked by adoption rate descending.

\begin{numbersbox}
\textbf{Key numbers generated:}\\[2pt]
\begin{tabular}{@{}ll@{}}
Repos with country (joined):      & 26,586 (44.2\% of corpus) \\
Repos without country:            & 33,543 (55.8\%) \\
Countries with $\geq$30 repos:    & \textbf{65} \\
Countries below threshold:        & 44 (excluded from ranking) \\[4pt]
Top country:   & Belgium \quad 89 repos \quad 89.9\% \\
Bottom country:& Rwanda \quad 49 repos \quad 22.4\% \\
Largest corpus:& United States \quad 3,608 repos \quad 82.5\% \\
Global (located): & 71.3\% \\
\end{tabular}
\end{numbersbox}

\textbf{Input:} \texttt{corpus\_detected\_full\_v2.csv} + \texttt{owner\_locations\_all.csv}

\textbf{Output:} \texttt{data/stage4/corpus\_with\_country.csv} (60,129 rows with
country column) and \texttt{data/stage4/country\_adoption.csv} (65-row ranking table)

\begin{outputbox}
\textbf{Thesis contribution.}
This script produces the 65-country ranking that is the centrepiece of the
geographic analysis. The range from 89.9\% (Belgium) to 22.4\% (Rwanda)
--- a 67.5~pp spread --- is the primary evidence for \textbf{RQ3}: adoption
rates vary significantly across countries. The ranking directly feeds the
choropleth world map, the top/bottom-10 bar charts, and the MENA deep-dive.
It also reveals the major corpus countries: the US (3,608 repos, 82.5\%),
China (4,196, 77.2\%), India (3,447, 51.2\%), and Germany (1,155, 78.2\%).
\end{outputbox}

\newpage

\subsection*{\texttt{stage4b\_statistical\_tests.py} --- Formal Statistical Tests (RQ1--RQ3)}

\begin{infobox}
\textbf{Purpose:} Runs the three formal statistical tests, one per research
question, and produces the numerical evidence for each RQ's conclusion.
\end{infobox}

\textbf{Method and tests:}

\textbf{RQ1 --- Binomial test.} One-sided binomial test against $H_0{:}\ p \leq 5\%$.
The 5\% null corresponds to incidental one-time agent use. Rejection confirms
adoption is statistically non-trivial in absolute terms.

\textbf{RQ2 --- McNemar's test.} Paired test comparing per-repository detection
outcomes under Pass~1-only vs.\ the full four-pass pipeline. Uses the
continuity-corrected test statistic. The discordant cell $b$ (missed by Pass~1,
detected by full pipeline) gives the undercounting delta.

\textbf{RQ3 --- Chi-squared + Fisher exact + Benjamini-Hochberg FDR.}
A chi-squared test of independence across all 65 countries tests whether
adoption rates differ globally. Followed by 65 pairwise Fisher's exact tests
(each country vs.\ rest-of-world); multiple comparisons corrected with
Benjamini-Hochberg FDR at $\alpha = 0.05$.

\begin{numbersbox}
\textbf{Key numbers generated:}\\[2pt]
\begin{tabular}{@{}p{3.5cm}p{8.5cm}@{}}
\toprule
\textbf{RQ1} & 73.0\% adoption (43,895/60,129). Wilson 95\% CI: [72.65\%, 73.35\%].
               Binomial $p < 10^{-300}$. \textbf{Reject $H_0$}. \\[3pt]
\textbf{RQ2} & Single-signal: 41.6\%. Full pipeline: 73.0\%. Delta: \textbf{+31.41 pp}.
               $b = 18{,}885$, $c = 0$. McNemar $\chi^2 = 18{,}883$, $p < 10^{-300}$.
               \textbf{Reject $H_0$}. \\[3pt]
\textbf{RQ3} & $\chi^2 = 2{,}234.55$, $df = 64$, $p < 10^{-300}$.
               41/65 countries significantly different from global mean
               after BH FDR correction. \textbf{Reject $H_0$}. \\
\bottomrule
\end{tabular}
\end{numbersbox}

\textbf{Input:} \texttt{data/stage4/corpus\_with\_country.csv} (60,129 rows)

\textbf{Output:} \texttt{data/stage4/statistical\_tests.txt} (full test output),
\texttt{data/stage4/pairwise\_fisher.csv} (per-country Fisher results with BH correction)

\begin{outputbox}
\textbf{Thesis contribution.}
This is the script that directly answers all three research questions with
formal statistical evidence. The results are the final answers reported in
the Results chapter and the Research Questions table in \texttt{results\_summary.tex}.
Notably, $c = 0$ in the McNemar table means there is \textbf{not a single
repository} that the full pipeline missed but Pass~1 found --- every
additional detection by the full pipeline is a genuine new discovery.
\end{outputbox}

\newpage

\subsection*{\texttt{stage4c\_breakdowns.py} --- Ecosystem, Language, and Tool Breakdowns}

\begin{infobox}
\textbf{Purpose:} Produces adoption breakdowns by ecosystem (package manager),
primary language (TypeScript vs.\ JavaScript), repository size, and AI coding
tool. Also analyses tool co-adoption (how many repos use multiple tools).
\end{infobox}

\textbf{Method.}
For each grouping variable, the script computes: total repos, agentic count,
ai-assisted count, adoption rate, Wilson~95\%~CI, and agentic-only rate.
For the tool breakdown, the script parses the \texttt{agents\_found} column
across all passes, splits multi-agent strings (e.g.\ \texttt{Claude Code|Cursor})
into individual tool counts, and counts how many unique tools each repository uses.
Tool pairs are generated by computing all unordered pairs from the multi-tool
string and counting pair frequencies.

\begin{numbersbox}
\textbf{Key numbers generated:}\\[2pt]
\begin{tabular}{@{}ll@{}}
Language breakdown: & TypeScript 77.1\% \quad JavaScript 60.5\% \\[2pt]
Ecosystem rates (top): & pnpm 82.5\% \quad deno 80.4\% \quad npm 73.7\% \quad vanilla 57.5\% \\[2pt]
Top tools (all repos): & Claude Code 26,854 (44.7\%) \quad Copilot 2,976 (4.9\%) \\
                       & Cursor 2,248 (3.7\%) \quad OpenAI Codex 2,127 (3.5\%) \\[2pt]
Tool loyalty:          & 1 tool: 40,532 (95.4\%)\quad 2 tools: 1,715 (4.0\%)\quad 3+: 176 (0.4\%) \\[2pt]
Top tool pair:         & Claude Code + Cursor (628 repos, 1.0\% of corpus) \\
\end{tabular}
\end{numbersbox}

\textbf{Input:} \texttt{data/stage4/corpus\_with\_country.csv}

\textbf{Output:} \texttt{breakdown\_ecosystem.csv}, \texttt{breakdown\_language.csv},
\texttt{breakdown\_size.csv}, \texttt{coadoption\_tool\_counts.csv},
\texttt{coadoption\_tool\_pairs.csv}, \texttt{coadoption\_per\_repo.csv}

\begin{outputbox}
\textbf{Thesis contribution.}
This script provides the data behind several key Discussion chapter arguments:
the TypeScript/JavaScript gap (77.1\% vs.\ 60.5\%) supports the argument
that typed codebases benefit more from AI suggestions due to richer context.
The pnpm lead (82.5\%) links engineering-focused tooling choices to AI adoption.
Claude Code's 44.7\% corpus share (59.2\% of adopters) establishes its
dominance in the npm ecosystem and motivates a discussion of why a
terminal-first CLI fits the npm development workflow. The 95.4\% single-tool
loyalty figure shows that despite many options, developers overwhelmingly
commit to one agent.
\end{outputbox}

\newpage

\subsection*{\texttt{stage4d\_size\_analysis.py} --- Repository Size and Adoption Rate}

\begin{infobox}
\textbf{Purpose:} Tests whether larger repositories have higher AI coding
agent adoption rates. Uses three independent operationalisations of
``size'' to rule out artefacts of any single metric.
\end{infobox}

\textbf{Method.}
Three size proxies are computed independently:
\begin{itemize}
  \item \textbf{Approach A --- Stars} (popularity proxy): fixed thresholds
        $\leq$50 (small), 51--500 (medium), $>$500 (large). Thresholds
        inherited from Stage~1a filters.
  \item \textbf{Approach B --- Commits} (development activity proxy):
        corpus tertiles. P33 = 114, P66 = 286 commits.
  \item \textbf{Approach C --- Code volume} (lang\_bytes proxy):
        corpus tertiles. P33 = 0.50~MB, P66 = 1.45~MB.
\end{itemize}
Approaches B and C use corpus tertiles so each stratum holds $\approx$33\%
of repositories, giving balanced groups. Approach A uses fixed thresholds
because the star distribution is extremely skewed (82.4\% have $\leq$50 stars).

\begin{numbersbox}
\textbf{Key numbers generated:}\\[2pt]
\begin{tabular}{@{}lrrrrr@{}}
\toprule
Approach & Small & n(small) & Medium & n(med) & Large \\
\midrule
A: Stars          & 70.5\% & 49,562 & 82.7\% & 8,375 & 91.7\% \\
B: Commits        & 64.8\% & 19,891 & 73.4\% & 19,863 & 80.6\% \\
C: Code volume    & 59.7\% & 19,843 & 73.3\% & 19,842 & 85.6\% \\
\bottomrule
\end{tabular}\\[4pt]
Range (large minus small): A: $+$21.2~pp \quad B: $+$15.8~pp \quad C: $+$26.0~pp
\end{numbersbox}

\textbf{Input:} \texttt{data/stage4/corpus\_with\_country.csv}

\textbf{Output:} \texttt{size\_stars.csv}, \texttt{size\_commits.csv},
\texttt{size\_langbytes.csv}, \texttt{size\_comparison.csv}

\begin{outputbox}
\textbf{Thesis contribution.}
All three independent operationalisations confirm the same monotonic pattern:
larger repositories adopt AI coding agents at higher rates. The consistency
across three different metrics rules out the possibility that the pattern
is an artefact of a particular size definition. The +26.0~pp range for code
volume (Approach~C) is the largest gap, suggesting that developers with
substantial codebases benefit most from --- and therefore are more likely
to adopt --- AI agents. This finding is discussed in the Discussion chapter
as evidence that adoption is driven partly by the perceived utility of agents
for complex, large-scale work.
\end{outputbox}

\newpage

\subsection*{\texttt{stage4e\_mena\_analysis.py} --- MENA Regional Deep-Dive}

\begin{infobox}
\textbf{Purpose:} Applies the same statistical analysis pipeline as the global
study to the MENA (Middle East and North Africa) region, enabling a focused
examination of within-region variation and a comparison of MENA against the
global benchmark.
\end{infobox}

\textbf{Method.}
The MENA region is defined as 19 countries across five sub-regions:
Gulf (UAE, Saudi Arabia, Kuwait, Qatar, Bahrain, Oman),
Levant (Israel, Jordan, Lebanon, Syria, Iraq),
North Africa (Egypt, Morocco, Algeria, Tunisia, Libya),
Turkey, and Iran. The same MIN\_REPOS = 30 threshold as the global analysis
is applied; 9 countries are eligible, 10 fall below the threshold.
Statistical tests: Fisher's exact (MENA aggregate vs.\ global 73.0\%),
chi-squared within MENA (9~$\times$~2 contingency table), and pairwise
Fisher's exact per country vs.\ MENA mean with BH FDR correction.

\begin{numbersbox}
\textbf{Key numbers generated:}\\[2pt]
\begin{tabular}{@{}ll@{}}
Total MENA repos (located):      & 1,446 (2.4\% of corpus, 5.4\% of located) \\
Eligible countries:              & 9 of 19 \\
MENA overall adoption:           & \textbf{65.1\%} [62.6\%, 67.5\%] \\
vs.\ global (Fisher):            & $p < 10^{-10}$, OR = 0.68 \\
Within-MENA ($\chi^2$):          & 147.52, $df=8$, $p \approx 10^{-28}$ \\
Countries significant (BH FDR):  & 7 of 9 \\[4pt]
Sub-region rates:                & Gulf 85.0\% \quad Levant 81.5\% \quad Turkey 70.7\% \\
                                 & North Africa 50.1\% \quad Iran 44.8\% \\[4pt]
Highest MENA country:            & Israel 85.6\% (n=181) \\
Lowest MENA country:             & Morocco 39.2\% (n=79) \\
Claude Code loyalty (MENA):      & 58.6\% (global: 59.2\%) \\
\end{tabular}
\end{numbersbox}

\textbf{Input:} \texttt{data/stage4/corpus\_with\_country.csv}

\textbf{Output:} \texttt{data/stage4/mena/} directory: \texttt{mena\_country\_ranking.csv},
\texttt{mena\_subregion\_breakdown.csv}, \texttt{mena\_statistical\_tests.txt},
\texttt{mena\_size\_breakdown.csv}, \texttt{mena\_tool\_breakdown.csv},
and additional CSV files for coverage and co-adoption.

\begin{outputbox}
\textbf{Thesis contribution.}
The MENA analysis demonstrates that the global geographic variation found in
RQ3 is replicated within a single cultural/geographic macro-region: a 46.4~pp
spread exists within MENA alone (Israel 85.6\% to Morocco 39.2\%).
The Gulf--Levant sub-regions (85.0\%, 81.5\%) outperform the global average
(73.0\%), while North Africa (50.1\%) and Iran (44.8\%) lag significantly
behind. Crucially, Claude Code loyalty is nearly identical inside MENA
(58.6\%) and globally (59.2\%), showing that tool preferences do not drive
the regional gap --- structural factors (internet access, GPU access,
economic constraints, language barriers) are more likely explanations,
discussed in the Discussion chapter.
\end{outputbox}

\end{document}
